\section{Introduction}

Nonabelian Hodge theory identifies two moduli spaces whose cohomology carries very different geometric filtrations. On the Higgs-bundle side, the Hitchin map defines the perverse Leray filtration. On the character-variety side, the mixed Hodge structure defines the weight filtration. The $P=W$ conjecture asserts that these filtrations coincide after doubling the indices of the first.

The geometric setting comes from Hitchin's moduli spaces and integrable systems \cite{Hitchin1987SelfDuality,Hitchin1987Integrable}. The analytic correspondence rests on the foundational work of Hitchin, Donaldson, Corlette, and Simpson \cite{Hitchin1987SelfDuality,Donaldson1987,Corlette1988,Simpson1992}. We use its twisted, fixed-determinant form, including the action by torsion line bundles, as described by Hausel and Thaddeus \cite[Sections~2 and~5]{HT2003}.

Let $C$ be a smooth projective complex curve of genus $g\ge2$. Fix an integer $n\ge2$, an integer $d$ coprime to $n$, and a line bundle $L\in\Pic^d(C)$. Write
\[
X=M_{\mathrm{Dol}}^L(\SL_n)
\]
for the moduli space of stable rank-$n$ Higgs bundles $(E,\theta)$ with $\det E\simeq L$ and $\tr\theta=0$. Here stability means that every proper nonzero $\theta$-invariant subbundle has slope smaller than $d/n$. Its Hitchin map is
\[
h:X\longrightarrow A=\bigoplus_{i=2}^nH^0(C,K_C^i),
\]
given by the coefficients of the characteristic polynomial. Put
\[
R=(n^2-1)(g-1),\qquad b=\frac{n(n-1)}2,
\qquad \zeta=\exp(2\pi\sqrt{-1}d/n).
\]
The corresponding twisted character variety is
\begin{equation}\label{intro:betti}
X^B=\left\{(A_1,B_1,\ldots,A_g,B_g)\in\SL_n(\C)^{2g}:
\prod_{j=1}^g[A_j,B_j]=\zeta I_n\right\}\mathbin{/\mkern-6mu/}\SL_n(\C),
\end{equation}
where the group acts by simultaneous conjugation and $[A,B]=ABA^{-1}B^{-1}$. Coprimality makes both spaces smooth and makes every representation in \eqref{intro:betti} irreducible. Nonabelian Hodge theory gives a diffeomorphism between them \cite{HT2003,Simpson1997}.

We use rational cohomology unless complex coefficients are indicated. The Hitchin map is proper, $\dim X=2R$, $\dim A=R$, and
\[
\dim(X\times_A X)-\dim X=R.
\]
Thus the standard $P=W$ normalization is
\begin{equation}\label{intro:P}
P_kH^m(X)=\im\left\{
H^{m-R}\!\left(A,{}^p\tau_{\le k}Rh_*\Q_X[R]\right)
\longrightarrow H^m(X)\right\}.
\end{equation}
Here ${}^p\tau$ denotes truncation for the middle perverse $t$-structure. We recall the dimension and shift conventions in Section~\ref{sec:lefschetz}. Transport $P$ to $H^*(X^B)$ by nonabelian Hodge theory, and let $W$ denote Deligne's weight filtration there \cite{Deligne1971,Deligne1974}.

\begin{theorem}\label{thm:main}
For every smooth projective complex curve $C$ of genus $g\ge2$, every composite integer $n\ge4$, every integer $d$ with $\gcd(n,d)=1$, and every $L\in\Pic^d(C)$, one has
\[
P_kH^m(X^B,\Q)=W_{2k}H^m(X^B,\Q)
=W_{2k+1}H^m(X^B,\Q)
\]
for all $m,k\ge0$.
\end{theorem}

The assertion concerns the entire cohomology. The group
$\Gamma=\Pic^0(C)[n]$ acts on $X$ by tensor product. Its averaging projector
\[
\operatorname{av}_{\Gamma}=\frac1{|\Gamma|}\sum_{\gamma\in\Gamma}\gamma^*
\]
has image $H^*(X,\Q)^\Gamma$ and kernel the \emph{variant cohomology}
$H^*_{\mathrm{var}}(X,\Q)$. Thus, after complexification, the theorem includes every nontrivial $\Gamma$-character, not only the invariant part. Combining Theorem~\ref{thm:main} with the established prime-rank cases \cite[Theorem~0.3]{MS2024} resolves the smooth coprime fixed-determinant, trace-free $P=W$ conjecture in every rank $n\ge2$.

\subsection{History and the fixed-determinant problem}

The weight symmetries behind $P=W$ arose from the arithmetic study of character varieties by Hausel and Rodriguez-Villegas \cite{HRV2008}, extended to generic puncture data by Hausel, Letellier, and Rodriguez-Villegas \cite{HLRV2011}. De Cataldo, Hausel, and Migliorini formulated the filtration comparison and established it in rank two \cite{dCHM2012}. Their theorem relates the topology of an algebraic integrable system to the mixed Hodge theory of an affine character variety: relative hard Lefschetz on the Higgs side explains curious hard Lefschetz on the Betti side.

Tautological classes have been central to the higher-rank problem. For $\GL_n$, generation by these classes was established in rank two by Hausel and Thaddeus \cite{HTGenerators2004} and in arbitrary rank by Markman \cite[Theorem~7]{Markman2002}; Shende determined their Betti weights \cite{Shende2017}. De Cataldo, Maulik, and Shen proved $P=W$ for $\GL_n$ in genus two and obtained the required perversity bounds for tautological classes in arbitrary genus \cite[Theorems~0.2 and~0.4--0.6]{dCMSJAMS2022}.

Maulik and Shen proved the coprime $\GL_n$ conjecture in arbitrary rank and genus \cite{MS2024}. Further proofs use the Hecke-algebra action of Hausel, Mellit, Minets, and Schiffmann \cite{HMMS2025} and the Fourier-transform method of Maulik, Shen, and Yin \cite[Theorem~0.4]{MSY2025}.

For fixed determinant, the quotient by $\Gamma$ detects the invariant cohomology, and the projective tautological generation theorem applies to this quotient \cite[Theorem~2.1(i)]{MS2024}. It does not generate the variant part. De Cataldo, Maulik, and Shen determined the two filtrations on the variant part in prime rank and reduced $\SL_p$ to $\GL_p$ \cite{dCMS2022}; together with the general $\GL_p$ theorem, this gives the prime-rank cases cited above. In arbitrary rank, Maulik and Shen constructed endoscopic correspondences compatible with the perverse filtrations \cite[Theorem~5.4]{MSEndoscopy2021}, without identifying the full filtrations in composite rank. Corollary~\ref{cor:endoscopic-weights} combines our all-rank equality with cover-side $\GL_r$ $P=W$ to give the exact Betti-weight image equality asked in \cite[Question~5.5]{MSEndoscopy2021}, for the canonical-twist operator transported by nonabelian Hodge theory.

Our argument constructs an operator on the full cohomology, before taking $\Gamma$-invariants. Its algebraic pattern follows the common $\mathfrak{sl}_2$ structure and flag-correspondence methods of Hausel, Mellit, Minets, and Schiffmann \cite[Proposition~7.9 and Section~8]{HMMS2025}. The fixed-determinant construction requires two further geometric steps. First, we pull Mellit's toric stratifications through a determinant-root cover and prove integral local constancy and ordinary base change. The inputs include the relative product stratification over ordered eigenvalues, including collisions \cite[Section~7.4, Theorem~7.5, and proof of Proposition~8.7]{Mellit2025}. This yields curious hard Lefschetz on the full fixed-determinant cohomology, rather than only its invariant part.

Second, we use a finite diagram of algebraic fibers to give weight zero to the cochains of a topological moving-point base. Shende used this weight convention for simplicial bases and tautological classes \cite[proof of the main Theorem]{Shende2017}. Here the construction also carries variations and proper Gysin maps, using mixed Hodge modules and their coherent enhancement \cite{SaitoExtension1990,TubachStacks2025}. Its finite weight-graded monodromy supplies the lowering operator. These diagram weights are distinct from the usual weights of an algebraic moving-point curve; that distinction is needed for integration over the curve to have weight shift zero.

\subsection{The operator behind the comparison}

The universal projective bundle on $C\times X$ has local vector-bundle lifts $\mathcal E$. The following characteristic class is independent of the lift and descends to $C\times X$; integration over $C$ gives
\[
e=\int_C\left(\ch_2(\mathcal E)-\frac{c_1(\mathcal E)^2}{2n}\right)
=\frac1{2n}\int_C\ch_2(\End\mathcal E)\in H^2(X).
\]
We also write $e$ for cup product by this class. Section~\ref{sec:lefschetz} proves that $H^2(X)=\Q e$ and that $e$ satisfies relative hard Lefschetz, with center $R$. Section~\ref{sec:betti} proves curious hard Lefschetz for the same class and the half-weight grading, again with center $R$.

It remains to construct one degree-$-2$ operator $f$ such that
\begin{equation}\label{intro:operator-goal}
fP_k\subset P_{k-2},\qquad
fW_j\subset W_{j-4},\qquad [[f,e],e]=-2e.
\end{equation}
An elementary Lefschetz-module argument then identifies both filtrations as sums of eigenspaces of the same operator $[e,f]$. It also handles extensions between filtration grades.

To construct $f$, allow one logarithmic pole, move it over a finite cover $S\to C$, and choose a full flag at the pole. The logarithmic $\lambda$-connection parameter interpolates between Higgs bundles at $\lambda=0$ and connections at $\lambda=1$. The underlying moduli constructions and semistable reduction come from \cite{dCFHLogHodge,FHZ2025}; the restriction argument uses the semiprojective geometry of \cite{HRV2015}. We verify the fixed determinant, residue, and flag conditions, as well as compatibility with the full nonabelian-Hodge identification. Write $X_\lambda$ for the fixed-pole scalar-residue moduli space, with $X_0=X$. At scalar residue, the flag space fits into a correspondence
\[
X_\lambda\xleftarrow{\ \pr\ }S\times X_\lambda
\xleftarrow{\ \pi\ }Z_\lambda
\xrightarrow{\ \iota\ }Y_{\lambda,0}.
\]
Here $Y_{\lambda,0}$ permits a flag-preserving residue with prescribed scalar diagonal, while $Z_\lambda$ requires the residue itself to be scalar. We first deform the ordered residues to distinct values and shift two residue eigenline lattices in opposite directions. Restriction isomorphisms from the total logarithmic Hodge family transport the resulting translation to a degree-zero operator $T$ on $H^*(Y_{\lambda,0})$. On the connection side, $T$ is eigenvalue monodromy. A power has unipotent action, and its logarithm $N$ lowers the diagram weights by two.

At $\lambda=0$, the resulting operator, before multiplication by a nonzero scalar, is
\[
f'=\pr_*\pi_*\iota^*N^2\iota_*\pi^*\pr^*.
\]
The flag shifts cancel. The square $N^2$ lowers weights by four, and integration over $S$ has weight shift zero in the diagram construction. On the Higgs side, nearby cycles realize this same operator as a morphism
$Rh_*\Q_X\to Rh_*\Q_X[-2]$, giving the perverse bound. A lattice Riemann--Roch calculation then gives the double commutator in \eqref{intro:operator-goal} after explicit rescaling.

The proof keeps three comparison statements separate: ordinary base change for the Betti family, restriction isomorphisms for the logarithmic Hodge family, and specialization over the Hitchin base. Each is proved with its own hypotheses. Section~\ref{sec:lefschetz} gives the Lefschetz comparison criterion. Section~\ref{sec:betti} proves ordinary base change and curious hard Lefschetz for the full fixed-determinant Betti family. Section~\ref{sec:diagram} constructs the diagram weights, and Section~\ref{sec:logarithmic} establishes the logarithmic restriction isomorphisms and full-cohomology comparison. Section~\ref{sec:translation} constructs the common operator and proves its weight bound; Section~\ref{sec:specialization} proves its perverse bound by specialization. Section~\ref{sec:commutator} computes the double commutator, and Section~\ref{sec:conclusion} assembles these results to prove Theorem~\ref{thm:main}.
